\section{Every Graph is the Limit of Graphs with at most $3$ Vertices}\label{appendix: tree shape important}

\begin{figure}
    \centering
    \tikzfig{graph_co}
    \caption{An illustration of the construction in Lemma~\ref{lemm:thin_build}}
    \label{fig:graph_co}
\end{figure}


\begin{lemma}\label{lemm:thin_build}
    Every graph is isomorphic to a limit of a graph co-decomposition whose bags and adhesions each contain at most three vertices.
\end{lemma}
\begin{proof}
    See Figure~\ref{fig:graph_co}.
    Let $G$ be a graph. Let $(x_1,\ldots,x_n)$ be an ordering of the vertex set of $G$. Since otherwise it is trivial, suppose that $n\ge 3$. Define the following partitions of $V(G)$:
    
    Let $\Pa_i=\{\;\{x_1,\ldots,x_{i-1}\},\;\{x_{i}\},\;\{x_{i+1},\ldots,x_n\}\;\}$ for each $i\in\{2,\ldots,n-1\}$,
    
    $\Pa_1=\{\;\{x_1\},\;\{x_{2},\;\ldots,x_n\}\;\}\quad$ and $\quad\Pa_n=\{\;\{x_1,\ldots,x_{n-1}\},\;\{x_{n}\}\;\}$.
    
    Let $\Ea_i=\{\;\{x_1,\ldots,x_{i}\},\;\{x_{i+1},\ldots,x_n\}\;\}$ for each $i=[n-1]$.
    
    Let $\Da_i=\{\;\{x_i\},\; V(G)\setminus\{x_{i}\}\;\}$ for each $i\in[n]$.
    
    Let $\Qa_{i,j}=\{\;\{x_i\},\;\{x_j\},\; V(G)\setminus\{x_i,x_j\}\;\}$ for each $i,j\in[n]$.
    
    \noindent For a partition $\Xa$ of $V(G)$, denote by $G/\Xa$ the graph with vertices the parts of~$\Xa$ and an edge between two parts (perhaps as a loop between the same two parts) when there is an edge in $G$ between a vertex in each of the parts. There is a homomorphism from $G$ to $G/\Xa$ taking each vertex to its part, which we call here the \define{projection to $G/\Xa$}.

    Let $H$ be the graph with the following vertices: $p_i$ for $i\in[n]$, $q_{i,j}$ for $i,j\in[n]$ and edges $e_i$ spanning $(p_i,p_{i+1})$ for $i\in[n-1]$ and $d_{i,j}^k$ for $k\in\{i,j\}$ and $i,j\in[n]$ spanning $(p_k,q_{i,j})$. For $i\in [n]$, let $G_{p_i}$ be a copy of $G/\Pa_i$, let $G_{q_{i,j}}$ be a copy of $G/\Qa_{i,j}$, let $G_{e_i}$ be a copy of $G/\Ea_i$ and $G_{d_{i,j}^k}$ be a copy of $G/\Da_k$. Let $\Ha:=(H,\{G_a\},\{g_t^e\})$ be the graph co-decomposition given by these definitions with maps defined so that a vertex $x$ in $G/\Xa$ is mapped to a vertex $y$ in $G/\Ya$ if the part of $\Xa$ corresponding to $x$ is a subset of the part of $\Ya$ corresponding to $y$. By definition these maps commute with the projections to each quotient, which we denote by $g_t:G\rightarrow G_t$ for each $t\in V(H)$. Observe that the pair $(G,\{g_t\})$ is a cone over $\Ha$. We will prove that it is a limit cone.

    Let $(G',\{f_t\})$ an arbitrary cone over $\Ha$. We are required to show that there is a unique homomorphism $\phi:G'\rightarrow G$ which commutes with the maps $\{g_t\}$ and $\{f_t\}$.

    \noindent\textbf{Claim.} For every $u\in V(G')$, there is a unique $i\in [n]$ so that the image of $u$ in $G_a$ is the partition containing $x_i$ for each $a\in V(H)\cup E(H)$.
    
    \noindent\textit{Proof of claim.} Let $j\in [n]$. If $u$ is mapped by $f_{p_j}$ to $\{x_1,\ldots,x_{j-1}\}$ in $G_{p_j}=\Pa_j$, then it is mapped to $\{x_1,\ldots,x_{j}\}$ in $G_{e_{j}}=\Ea_j$ and hence to $\{x_1,\ldots,x_{j+1}\}$ in $G_{p_{j+1}}=\Pa_{j+1}$. If $u$ is mapped instead by $f_{p_j}$ to $\{x_{j}\}$ in $G_{p_j}=\Pa_j$, then it is mapped to $\{x_1,\ldots,x_{j}\}$ in $G_{e_{j}}=\Ea_j$ and $\{x_j,\ldots,x_n\}$ in $G_{e_{j-1}}=\Ea_{j-1}$ and hence to $\{x_1,\ldots,x_{j+1}\}$ in $G_{p_{j+1}}=\Pa_{j+1}$ and $\{x_j,\ldots,x_{n}\}$ in $G_{p_{j-1}}=\Pa_{j-1}$. Finally, if $u$ is mapped by $f_{p_j}$ to $\{x_{j+1},\ldots,x_n\}$ in $G_{p_j}=\Pa_j$, then it is mapped to $\{x_j,\ldots,x_n\}$ in $G_{e_{j-1}}=\Ea_{j-1}$ and hence to $\{x_j,\ldots,x_n\}$ in $G_{p_{j-1}}=\Pa_{j-1}$.

    We conclude that there is a unique $i\in[n]$ so that $u$ is mapped by $g_{p_i}$ to the singleton $\{x_i\}$ in $G_{p_i}=\Pa_i$ and for $k<i$, to the initial segment in $G_{p_k}=\Pa_k$ and for~$k>i$ to the final segment in $G_{p_k}=\Pa_k$.
    
    It remains to check the bags with indices $q_{k,j}$. This is routine after observing that~$u$ must map to the singleton $\{x_i\}$ in $G_{d_{i,j}^i}$ and the complement of $\{x_j\}$ in $G_{d_{i,j}^j}$ for each~$j\neq i$. This completes the proof of our claim. \hfill$\lozenge$

    Let $\phi:G'\rightarrow G$ be the map sending each vertex $u$ to the vertex $x_i$ for the $i$ given by the above claim. By construction, when we restrict only to vertices, this map commutes with the cones. Furthermore, it is easy to see that it is the unique map with this property. It remains to check $\phi$ defines a graph homomorphism. Suppose there is an edge of $G'$ which spans the vertices $u$ and $v$. Then there are indices~$i$ and~$j$ so that $\phi(u)=x_i$ and $\phi(v)=x_j$. Now consider the map $f_{q_{i,j}}:G'\rightarrow G_{q_{i,j}}$. Since this is a homomorphism, there must be an edge between the singletons $\{x_i\}$ and $\{x_j\}$ in~$G_{q_{i,j}}$. By definition, this only occurs when there is an edge from $x_i$ to $x_j$ in $G$. This completes our proof. 
\end{proof}

\begin{example}
    The `three' in Lemma~\ref{lemm:thin_build} is tight. Indeed, the graph co-decompositions with bags and adhesion of size at most two have limits which are direct products of graphs with at most two vertices each. Basic graphs such as the triangle are not realised as such.
\end{example}